\section{Analisi}
\subsection{Descrizione del Sito}
Il progetto nasce dall’esigenza di offrire un servizio farmaceutico moderno, accessibile e vicino alle esigenze di tutta la comunità. La farmacia si presenta come un punto di riferimento che coniuga la tradizione con l’innovazione, grazie alla realizzazione di un portale online all’avanguardia.

L’obiettivo principale è mettere la tecnologia al servizio della salute, garantendo un’esperienza digitale intuitiva per i più giovani e strumenti semplici e chiari per gli utenti meno esperti. Il sito è stato progettato per essere facilmente fruibile da persone di tutte le età, con particolare attenzione all’accessibilità: sono stati adottati standard che permettono la navigazione anche a chi presenta disabilità visive, motorie o cognitive.

Attraverso questa piattaforma, la farmacia vuole abbattere le barriere fisiche e digitali, offrendo servizi personalizzati, informazioni chiare e la possibilità di interagire direttamente con i professionisti del settore. In questo modo, si rafforza il legame con la comunità, promuovendo la salute come diritto fondamentale per tutti.

\subsection{Analisi del utenza}

L’utenza di una farmacia è estremamente eterogenea e comprende persone di tutte le età, condizioni sociali e livelli di alfabetizzazione digitale. Un’attenzione particolare va rivolta alle persone affette da disabilità, che rappresentano una parte significativa della popolazione e che spesso incontrano ostacoli nell’accesso ai servizi sanitari, sia fisici che digitali.

Le principali categorie di disabilità da considerare sono:
\begin{itemize}
    \item \textbf{Disabilità visive}: includono ipovisione, cecità totale o parziale, daltonismo. Questi utenti necessitano di interfacce compatibili con screen reader, testi ad alto contrasto e possibilità di ingrandire i contenuti.
    \item \textbf{Disabilità motorie}: possono limitare l’uso di mouse o tastiera. È fondamentale garantire la navigazione tramite tastiera e prevedere pulsanti e link facilmente selezionabili.
    \item \textbf{Disabilità cognitive}: includono difficoltà di apprendimento, memoria o attenzione. L’interfaccia deve essere semplice, con istruzioni chiare e percorsi guidati.
\end{itemize}

Per rispondere alle esigenze di queste persone, il sito della farmacia è stato progettato seguendo le linee guida internazionali sull’accessibilità (WCAG), adottando soluzioni tecniche e grafiche che favoriscono l’inclusione. In questo modo, si garantisce a tutti gli utenti la possibilità di accedere ai servizi, alle informazioni sui farmaci e di interagire con i professionisti, riducendo il rischio di esclusione e promuovendo una reale equità nell’accesso alla salute.